\chapter{Ant or other insect}
The \textbf{ant or other insect} is a monster class that appears in \textit{NetHack}, and is represented by the lowercase a glyph (\whitesym[a]). Ants and insects are designated internally by the macro \textbf{S_ANT}.

The monster class contains the following monsters:
\begin{myitemize}
\item{}\monsymlink[giant ant]
\item{}\monsymlink[killer bee]
\item{}\monsymlink[soldier ant]
\item{}\monsymlink[fire ant]
\item{}\monsymlink[giant beetle]
\item{}\monsymlink[queen bee]
\end{myitemize}

\section{Common traits}
Ants and insects are are neutral-aligned animals that have limbs, but no hands. They are carnivorous with the exception of both types of bee, and are also oviparous with the exception of the giant beetle. Ants and other insects attack with bites or stings, some of which can inflict strength-draining poison, and mostly lack any MR score (with the fire ant being the only exception at a meager 10 MR). Their corpses grant very little nutrition, and while most of them can provide useful poison resistance, those corpses will also be poisonous themselves.

Centipedes and spiders are grouped into the arachnid or centipede monster class, \whitesym[s], as they are arthropods but not considered insects (though NetHack is not real life, so do not always take this categorization for granted). Wholly-fantastic insects such as the xan constitute their own monster class.

\subsection{Generation}
Randomly generated ants and insects monsters are always created hostile.

The clerical monster spell summon insects will generate ants and insects unless the monster class is extinct or wiped out by genocide, in which case snakes are generated instead.

The ant or other insect is the first quest monster class for Valkyries, and makes up 24⁄175 of the monsters that are randomly generated on the Valkyrie quest. Random ants and other insects may also be generated on floors below the home level of this quest branch at level creation: one is generated on each of the upper filler level, the lower filler level(s) and the locate level, and two are generated on the goal level.

\section{Strategy}
Most of the insect monsters (particularly the ants) have high speed along with attacks that are decently strong and can sometimes inflict poison, making them early-game menaces that are overall tough to contend with. The soldier ant in particular is infamous for being a frequent source of game-ending stings even with more experienced players, and are dreaded encounters if poison resistance is not obtained by the time they appear: even with poison resistance, their attack frequency can easily cut into a hero's HP, and their AC often makes them hard to kill in return. The giant beetle is the slowest of the group in contrast, with a movement speed of 6 and the weakest AC at 4, and it cannot inflict poison through its attacks—however, their bite is the strongest of the monster class and makes melee combat similarly risky.

Even if a hero is not outright killed from severe poison damage, the attribute loss usually warrants finding a cure as soon as possible: successful prayer, potions of restore ability, and the restore ability spell are methods that can restore attributes lost to poison. This also makes it imperative to find and identify a poison resistance source as well, since while one poisoned sting from an insect monster will not always put the hero under, the HP damage and possible attribute loss from multiple stings back to back is liable to leave a hero that survives in dire straits.

Elbereth is useful in driving off aggressive insects, and their low MR score makes them very vulnerable to spells, wands and potions. Chain lightning is an ideal spell for handling insect swarms. While genocide is usually unnecessary to deal with any of the monsters, ants and insects are usually a valid "emergency" target if they have the hero trapped with few other means of escaping or defeating them. A tinning kit can be used on poisonous corpses in order to make them safe to eat and potentially gain poison resistance. Once a hero gains poison resistance along with enough quality equipment and/or weapon skill to reliably dispatch ants and bees, they will usually not be a problem for the rest of the game unless the hero becomes careless.

The summon insects spell is typically not seen until near the end stages of the game (e.g. hostile priests in Moloch's Sanctum and on the Astral Plane), when the insects generated by the spell are generally a nuisance at worst that get between you and attacking the hostile caster in question. Incidentally, this can also serve as a limited form of "cushion" between you and much stronger hostiles, with the exception of the Riders: due to slower HP regeneration at higher experience levels it may be best not to idle (even with the regeneration property), unless you have a method of healing yourself thoroughly and need the time and space to apply it reliably. Similarly, they can also provide a vector for chain lightning spells, although the amount of hostiles a hero will be facing in late-game situations where insects are summoned means that significant power drain will usually occur—heroes conserving their energy for other spells should treat this as a form of last resort.

\section{History}
The giant ant and killer bee first appear in Hack 1.21, a port of Jay Fenlason's Hack, while the giant beetle appears in Hack for PDP-11, and all three monsters are included in the initial bestiary for Hack 1.0. From these versions to NetHack 2.3e, the giant ant uses the \whitesym[A] glyph and is more similar in behavior to the modern soldier ant, while the killer bee uses the \whitesym[k] glyph—the giant beetle uses the \whitesym[b] glyph.

NetHack 3.0.0 establishes the ant or other insect monster class and moves the giant ant, killer bee and giant beetle to that class, alongside the newly-added soldier ant, fire ant and queen bee.

From NetHack 3.0.0 to NetHack 3.6.7, including some variants based on those versions, the killer bee and soldier ant have a difficulty of 5 and 6 respectively. Poison stings and other attacks can cause instant death in these versions, and applying a unicorn horn can cure strength drain from the attacks or the consumption of poisonous corpses. NetHack 5.0.0 changes the difficulty of the killer bee and soldier ant via commit 47724f01, removes the ability of poison attacks to cause instant death via commit d0b11fd2, and additionally removes the unicorn horn's ability to restore attributes via commit 43d331c4. Furthermore, homemade tins of insect meat grant 40 nutrition rather than the standard 5—this is changed in NetHack 5.0.0 via commit 39bd259b.

